\documentclass[10pt,a4paper]{amsart}
\usepackage[margin=19mm]{geometry}
\usepackage{amsmath,amssymb,mathtools}
\usepackage[scr=rsfs,cal=euler]{mathalfa}
\usepackage{xfrac}
\usepackage[unicode,hidelinks,bookmarks=false]{hyperref}
\newcommand{\ie}{i.e.\ }
\newcommand{\cf}{cf.\ }
\newcommand{\resp}{respectively\ }
\newcommand{\Subsection}{\subsection}
\providecommand{\nobreakdash}{\nobreak\hbox{-}}
\setlength{\emergencystretch}{2em}
\multlinegap0pt
\usepackage{bm}
\usepackage{bbm}
\usepackage{dsfont}
\usepackage{xspace}
\usepackage{cancel}
\usepackage{mathtools}
\usepackage{amsthm}
\makeatletter
\@ifundefined{theorem}{\newtheorem{theorem}{Theorem}}{}
\@ifundefined{definition}{\newtheorem{definition}[theorem]{Definition}}{}
\makeatother
\title[Self-graphing equations]{Self-graphing equations}
\author{Samuel Allen Alexander}
\date{Source: arXiv:2504.00006v1 (CC BY 4.0)}
\begin{document}
\maketitle

\noindent\emph{Source and licence.} Self-graphing equations. Version arXiv:2504.00006v1 (CC BY 4.0), distributed under the Creative Commons Attribution 4.0 International licence, as recorded in the arXiv licence metadata for this exact version and shown on the abstract page https://arxiv.org/abs/2504.00006v1. Full rights notice: licenses/arxiv-2504.00006-CC-BY-4.0.txt.\par\smallskip

\noindent\textit{Editorial scope.} Counted unit: the Theorem whose source label is \texttt{maintheorem} in the article (source0.tex lines 447--451, source label \texttt{maintheorem}), delivered together with the article's own complete original proof of it (source0.tex lines 453--504, one \texttt{proof} environment of 52 lines). No mathematics is written, repaired or re-derived: statement and proof are the article's own text line for line. Required context retained uncounted: recursionthm (lines 360--365); selfconstraintdefn (lines 389--401). Every definition, display and label of those lines is retained unchanged. Omitted from this package and not redistributed: 1--359, 366--388, 402--446, 452--452, 505--715. Nothing inside a retained excerpt is omitted or altered: every retained body is delivered exactly as the article's lines are written, so no omission receipt of any kind accompanies this package. Citations: the retained excerpts carry no citation command at all, so no bibliography entry of the article is transcribed and no reference is redistributed. Reference adaptation: none was needed and none was made; every reference inside the delivered unit resolves inside it, so no unresolved reference or citation is produced. Printed slips kept literal: no printed word, symbol, hypothesis or number of the counted statement or of its proof was altered. Layout and class: the article's own class is \texttt{report} with packages including amsmath, amsthm, latexsym, amssymb, graphicx, graphics, bold-extra, multirow, pdfpages, array, float, subcaption, bm, changepage; the wrapper is the lane's pinned \texttt{amsart} editorial wrapper at 10pt on A4, which restates the theorem-like environments the retained text needs and adds the article's own macro definitions that the retained text needs (none), with extra wrapper packages (bm, bbm, dsfont, xspace, cancel, mathtools). The delivered numbering is the unit's own. Assets: no retained span contains a figure environment, a graphics-inclusion command, a PDF-page inclusion, a table or a TikZ picture; no image file is copied into this package, the package declares no asset dependency and no image file is redistributed with it. Licence: version arXiv:2504.00006v1 of this article is distributed under the Creative Commons Attribution 4.0 International licence, as recorded in the arXiv licence metadata for this exact version and shown on the abstract page https://arxiv.org/abs/2504.00006v1. A full rights notice is delivered at licenses/arxiv-2504.00006-CC-BY-4.0.txt. Characters: every retained excerpt is pure ASCII, so the delivered engine, XeLaTeX with its default Latin Modern text font, reports no missing character.
\begin{theorem}
\label{recursionthm}
  (The Recursion Theorem)
  For every total computable $f:\mathbb N\to\mathbb N$, there is some
  $n\in\mathbb N$ such that $\varphi_n=\varphi_{f(n)}$.
\end{theorem}

\begin{definition}
\label{selfconstraintdefn}
  The glyphed notion of equations $\mathcal A$ is \emph{self-constrained}
  if there exists a G\"odel numbering $\ulcorner\bullet\urcorner$
  of $A^*$ and a total computable $f:\mathbb N\to\mathbb N$ such that:
  \begin{itemize}
    \item
    For all $n\in\mathbb N$, if $\varphi_n(0)=\ulcorner\tau\urcorner$
    for some $\tau\in A^*$, then
    $f(n)=\ulcorner\sigma\urcorner$ for some $\sigma\in A^*$
    such that $\mathrm{Gr}(\sigma)=\mathrm{Gl}(\tau)$.
  \end{itemize}
\end{definition}

\begin{theorem}
\label{maintheorem}
  If $\mathcal A$ is self-constrained then there exists a self-graphing
  equation in $\mathcal A$.
\end{theorem}

\begin{proof}
  Let $\ulcorner\bullet\urcorner$ and $f:\mathbb N\to\mathbb N$ be as in
  Definition \ref{selfconstraintdefn}.

  Subclaim: There is a total computable function $g:\mathbb N\to\mathbb N$
  such that for all $n\in\mathbb N$, $\varphi_{g(n)}(0)=f(n)$.

  This Subclaim is actually a special case of a theorem
  from computability theory called
  the ``$Smn$ theorem'', but we will sketch a direct proof here.
  Let $g:\mathbb N\to\mathbb N$ be the function computed by the following
  algorithm:
  \begin{enumerate}
    \item Take input $n\in\mathbb N$.
    \item Let $X=f(n)$.
    \item Let $P$ be the following algorithm:
    \begin{enumerate}
      \item Take input $m\in\mathbb N$.
      \item Output $X$ (ignoring the value of $m$).
    \end{enumerate}
    \item Output an encoding of $P$ (a number $k$ such that
      $\varphi_k$ is the function computed by $P$).
  \end{enumerate}
  For any $n\in\mathbb N$, by construction $\varphi_{g(n)}$ is the
  function computed by the above algorithm $P$ (for the given $n$).
  Thus $\varphi_{g(n)}(0)$ is computed by ignoring the input $m=0$
  and outputting $X=f(n)$. Thus $\varphi_{g(n)}(0)=f(n)$.
  Since we have provided an algorithm for $g$, $g$ is computable.
  Clearly $\mathrm{dom}(g)=\mathbb N$, so $g$ is total computable.
  This proves the Subclaim.

  Let $g$ be as in the Subclaim.
  By the Recursion Theorem (Theorem \ref{recursionthm}) there is some
  $n\in\mathbb N$ such that $\varphi_n=\varphi_{g(n)}$.
  In particular,
  \[
    \varphi_n(0) = \varphi_{g(n)}(0) = f(n) \mbox{ is defined. ($*$)}
  \]
  Let $\sigma,\tau\in A^*$ be such that $f(n)=\ulcorner\sigma\urcorner$
  and $\varphi_n(0)=\ulcorner\tau\urcorner$.
  We claim $\sigma$ is a self-graphing equation in $\mathcal A$.
  To see this,
  compute:
  \begin{align*}
    \mathrm{Gr}(\sigma)
      &= \mathrm{Gl}(\tau)
        &\mbox{(Definition \ref{selfconstraintdefn})}\\
      &= \mathrm{Gl}(\sigma),
        &\mbox{(By $*$, $\sigma=\tau$)}
  \end{align*}
  as desired.
\end{proof}

\end{document}
